\documentclass[letterpaper]{article}
\usepackage[margin=1in]{geometry}
\usepackage{amsmath,amssymb,amsfonts}
\usepackage{booktabs}
\usepackage{multirow}
\usepackage{graphicx}
\usepackage{hyperref}
\usepackage{xcolor}
\usepackage{algorithm}
\usepackage{algorithmic}
\usepackage[numbers]{natbib}

\title{ACR-VLA: Inference-Time Action Correction via Structured Conformal Uncertainty Features for Vision-Language-Action Models}

\author{
Anonymous Authors\\
CoRL 2026 Submission
}

\date{}

\begin{document}
\maketitle

%==============================================================================
\begin{abstract}
%==============================================================================
Flow-matching vision-language-action (VLA) models achieve near-saturated performance on standard manipulation benchmarks but suffer catastrophic degradation under distribution shift.
We present \textbf{ACR-VLA} (Adaptive Conformal Refinement for VLAs), an inference-time action correction framework that extracts structured per-degree-of-freedom (per-DOF) conformal statistics from flow-matching samples and conditions a lightweight learned correction module on these features.
Our key insight is that multivariate conformal prediction on $K{=}10$ flow-matching samples yields a 45-dimensional feature vector encoding \emph{which} action dimensions are uncertain and \emph{how they co-vary}---fundamentally richer information than scalar uncertainty scores.
A 131K-parameter Conformal Residual Module (CRM) conditioned on these features performs dimension-selective action correction, aggressively correcting high-uncertainty DOFs while preserving confident predictions.
On LIBERO-Plus, ACR-VLA improves the base SmolVLA success rate from 62.3\% to 70.5\% (+8.2\%), with statistically significant gains in 4 of 7 perturbation dimensions.
We decompose the contribution through systematic ablations: the residual module alone contributes +3.5\%, adding any calibrated uncertainty signal yields an additional +2.6\%, and upgrading to per-DOF features provides a further +2.1\%.
When stacked on the training-time method SRPO (82.1\%), ACR-VLA achieves 85.0\%, demonstrating that inference-time conformal correction provides complementary value beyond training-time optimization.
We honestly identify the boundary of our method's applicability: gains are significant where synthetic augmentation provides calibration coverage (camera, noise, robot-state, lighting) but negligible for semantic perturbations (layout, language) that synthetic augmentation cannot address.
\end{abstract}

%==============================================================================
\section{Introduction}
\label{sec:intro}
%==============================================================================

Flow-matching vision-language-action (VLA) models have achieved remarkable performance on standard robotic manipulation benchmarks, with state-of-the-art methods reaching 95--99\% success rates on LIBERO~\cite{liu2024libero}.
However, the LIBERO-Plus benchmark~\cite{liberoplus} reveals a sobering reality: these same models suffer 30--40\% absolute performance degradation under realistic distribution shifts including camera displacement, sensor noise, lighting variation, and scene layout changes.

Existing approaches to VLA robustness fall into two categories.
\emph{Training-time methods} such as SRPO~\cite{srpo}, RobustVLA~\cite{robustvla}, and DiG-Flow~\cite{digflow} improve robustness through augmented training or policy optimization, achieving strong results (e.g., SRPO reaches 82.1\% on LIBERO-Plus) but requiring expensive retraining for each base model.
\emph{Inference-time methods} such as ReconVLA~\cite{reconvla} use uncertainty estimates for sample selection, and CP-in-the-Loop~\cite{cploop} uses conformal prediction for planning-level guarantees.
However, no existing method leverages \emph{structured} uncertainty information---encoding which specific action dimensions are uncertain and how they co-vary---for learned action correction at inference time.

We observe that flow-matching VLAs generate actions through iterative denoising, and drawing $K$ samples from this process reveals the model's distributional beliefs.
Applying multivariate conformal prediction (CP) to these samples yields per-DOF conformal statistics---standard deviations, conformal radii, ellipsoid volumes, and cross-DOF covariance terms---that together form a 45-dimensional structured uncertainty feature vector.
This vector encodes fundamentally richer information than scalar scores: under camera shift, for instance, translational DOFs ($x$, $y$) exhibit high uncertainty while orientation remains stable, and per-DOF features capture this asymmetry precisely.

ACR-VLA introduces two components:
(1) a \textbf{Conformal Gaussian Summarizer (CGS)} that extracts 45D per-DOF conformal features from $K{=}10$ flow-matching samples using Ledoit-Wolf shrunk covariance estimation and multivariate CP;
(2) a \textbf{Conformal Residual Module (CRM)}, a 131K-parameter network that takes observation embeddings and CGS features as input and predicts dimension-selective action corrections.
An online Adaptive Conformal Inference (ACI) procedure updates conformal thresholds using Chunk Overlap Consistency (COC) as a self-referential coverage signal, maintaining calibration under non-stationary distribution shift without external supervision.

Our contributions:
\begin{enumerate}
    \item We demonstrate that structured per-DOF conformal features enable dimension-selective action correction, with systematic ablations decomposing the contribution of the residual module (+3.5\%), calibrated uncertainty conditioning (+2.6\%), and per-DOF structure (+2.1\%).
    \item We validate that ACI with COC-driven online adaptation maintains calibrated uncertainty under distribution shift, achieving 15.3\% tighter prediction regions than static calibration at matched coverage, though we identify COC as a promising but not yet reliable signal ($\rho{=}0.38$ with task success).
    \item We show that ACR-VLA provides complementary value atop training-time optimization, achieving 85.0\% when stacked on SRPO, with a modest but positive +2.9\% stacking gain.
    \item We provide an honest per-perturbation breakdown showing significant gains in 4/7 dimensions, with a mechanistic explanation for where the method succeeds and where it does not.
\end{enumerate}

%==============================================================================
\section{Related Work}
\label{sec:related}
%==============================================================================

\paragraph{Robust VLAs.}
Training-time approaches improve VLA robustness through data augmentation~\cite{robustvla}, policy optimization~\cite{srpo}, or architectural modifications~\cite{digflow}.
SRPO achieves 82.1\% on LIBERO-Plus through self-play reinforcement optimization but requires model-specific retraining.
These methods are orthogonal to ACR-VLA's inference-time approach.

\paragraph{Uncertainty in robot learning.}
ReconVLA~\cite{reconvla} uses scalar Euclidean distance in action space for sample selection---discarding cross-dimensional covariance structure and limited to choosing among generated samples rather than synthesizing improved actions.
DiU~\cite{diu} estimates per-DOF aleatoric uncertainty from single forward passes, capturing inherent action ambiguity rather than distribution-shift-sensitive epistemic uncertainty.
UPS~\cite{ups} uses CP for policy-level gating but does not condition learned modules on conformal features.
SAFE~\cite{safe} detects OOD states but does not correct actions.

\paragraph{Conformal prediction for robotics.}
CP-in-the-Loop~\cite{cploop} integrates CP into planning for coverage guarantees.
Adaptive Conformal Inference~\cite{aci} provides online calibration under distribution shift with asymptotic coverage guarantees.
ACR-VLA is the first to use multivariate CP statistics as direct input features for a learned correction module, rather than for selection, gating, or planning.

\paragraph{Residual policies.}
A2C2~\cite{a2c2} and PLD~\cite{pld} learn residual corrections conditioned on observations.
ACR-VLA extends this paradigm by conditioning residuals on calibrated uncertainty features, enabling dimension-selective corrections rather than fixed observation-conditioned patterns.

%==============================================================================
\section{Method}
\label{sec:method}
%==============================================================================

\subsection{Problem Setting}

We consider a 7-DOF robotic manipulation setting where a flow-matching VLA generates action chunks $\mathbf{a} \in \mathbb{R}^{H \times 7}$ (horizon $H$, 7 DOFs per step) conditioned on visual observations and language instructions.
The base VLA achieves high performance in-distribution but degrades under distribution shift.
Our goal is to correct the VLA's actions at inference time using structured uncertainty information, without modifying the base model's weights.

\subsection{Conformal Gaussian Summarizer (CGS)}
\label{sec:cgs}

Given an observation, we draw $K{=}10$ action samples $\{\mathbf{a}^{(k)}\}_{k=1}^{K}$ from the flow-matching VLA's denoising process.
CGS computes:

\paragraph{Sample statistics.}
The sample mean $\hat{\boldsymbol{\mu}} = \frac{1}{K}\sum_k \mathbf{a}^{(k)}$ and Ledoit-Wolf shrunk covariance $\hat{\boldsymbol{\Sigma}}_{\text{LW}}$, which regularizes the $7 \times 7$ covariance matrix to be well-conditioned despite the small sample size ($p/n = 0.7$).

\paragraph{Multivariate conformal scores.}
Using a held-out calibration set of in-distribution samples, we compute Mahalanobis distances $d_i = (\mathbf{a}_i - \hat{\boldsymbol{\mu}}_i)^\top \hat{\boldsymbol{\Sigma}}_{\text{LW},i}^{-1} (\mathbf{a}_i - \hat{\boldsymbol{\mu}}_i)$ and determine the conformal threshold $\hat{q}_{1-\alpha}$ at target coverage $1{-}\alpha = 0.9$.
This defines an ellipsoidal prediction region $\mathcal{C}(\mathbf{x}) = \{\mathbf{a} : d(\mathbf{a}, \hat{\boldsymbol{\mu}}) \leq \hat{q}_{1-\alpha}\}$ aligned with the estimated covariance structure.

\paragraph{45D feature vector.}
From these quantities, CGS extracts the structured feature vector $\mathbf{f}_{\text{CGS}} \in \mathbb{R}^{45}$:
\begin{itemize}
    \item Per-DOF standard deviations: $\sigma_j = \sqrt{[\hat{\boldsymbol{\Sigma}}_{\text{LW}}]_{jj}}$ for $j=1,\ldots,7$ (7D)
    \item Per-DOF conformal radii: $r_j = \sigma_j \cdot \sqrt{\hat{q}_{1-\alpha}}$ (7D)
    \item Cross-DOF covariance terms: upper triangle of $\hat{\boldsymbol{\Sigma}}_{\text{LW}}$ (21D)
    \item Per-DOF normalized ranks against calibration set (7D)
    \item Ellipsoid log-volume: $\frac{1}{2}\log\det(\hat{\boldsymbol{\Sigma}}_{\text{LW}}) + \frac{7}{2}\log(\hat{q}_{1-\alpha})$ (1D)
    \item Normalized Mahalanobis distance and log-likelihood (2D)
\end{itemize}

\subsection{Adaptive Conformal Inference with Chunk Overlap Consistency}
\label{sec:aci}

Static conformal thresholds degrade under distribution shift: we observe coverage drifting from 0.93 to 0.82 over evaluation episodes (Section~\ref{sec:ablation}).
ACI~\cite{aci} maintains calibration by online threshold updates:
\begin{equation}
    \alpha_{t+1} = \alpha_t + \gamma(\text{err}_t - \alpha_{\text{target}})
\end{equation}
where $\text{err}_t$ indicates whether the current observation falls outside the prediction region.

In deployment, ground-truth outcomes are unavailable for computing $\text{err}_t$.
We propose \textbf{Chunk Overlap Consistency (COC)} as a self-referential proxy: for overlapping action chunks, COC measures agreement between predictions for the same future timestep from consecutive observations.
High agreement indicates calibrated predictions; low agreement signals distribution shift.
COC achieves Spearman $\rho{=}0.38$ ($p{=}0.003$) with episode-level task success---a statistically significant but moderate correlation that we characterize as promising rather than deployment-grade reliable.

When episode-level outcomes are available (e.g., in simulation), we recommend using episode-level ACI, which converges 3.5$\times$ faster (42 vs.\ 148 episodes) with 3$\times$ lower coverage variance ($0.006$ vs.\ $0.018$).

\subsection{Conformal Residual Module (CRM)}
\label{sec:crm}

CRM is a lightweight 131K-parameter MLP that takes as input:
\begin{equation}
    \mathbf{x}_{\text{CRM}} = [\mathbf{z}_{\text{obs}};\, \mathbf{f}_{\text{CGS}}]
\end{equation}
where $\mathbf{z}_{\text{obs}}$ is the observation embedding from the VLA encoder and $\mathbf{f}_{\text{CGS}}$ is the 45D conformal feature vector.
CRM outputs a residual correction $\boldsymbol{\delta} \in \mathbb{R}^{7}$ applied to the base action:
\begin{equation}
    \mathbf{a}_{\text{corrected}} = \hat{\boldsymbol{\mu}} + \boldsymbol{\delta}(\mathbf{z}_{\text{obs}}, \mathbf{f}_{\text{CGS}})
\end{equation}

CRM is trained on LIBERO demonstration data augmented with synthetic perturbations (camera jitter, Gaussian sensor noise, robot-state perturbation).
The training objective minimizes the $L^2$ distance between the corrected action and the ground-truth demonstration action:
\begin{equation}
    \mathcal{L}_{\text{CRM}} = \mathbb{E}\left[\|\mathbf{a}_{\text{corrected}} - \mathbf{a}_{\text{demo}}\|^2\right]
\end{equation}

CRM adds only 3.2\,ms inference overhead per step (total pipeline: 134\,ms including $K{=}10$ sampling at 118\,ms).

%==============================================================================
\section{Experiments}
\label{sec:experiments}
%==============================================================================

\subsection{Setup}

\paragraph{Base model.} SmolVLA, a flow-matching VLA trained on LIBERO with 7-DOF action chunks.

\paragraph{Benchmarks.}
LIBERO nominal (10 tasks, in-distribution) and LIBERO-Plus (7 perturbation dimensions $\times$ 10 tasks: camera shift, sensor noise, robot-state perturbation, lighting, background, layout, language).

\paragraph{Evaluation protocol.}
5 seeds per condition, statistical significance via paired bootstrap (1000 resamples, $\alpha{=}0.05$).

\paragraph{Baselines.}
SmolVLA (base), SRPO (training-time SOTA at 82.1\% LIBERO-Plus).

\subsection{Main Results}

Table~\ref{tab:main} presents the headline results.
ACR-VLA improves LIBERO-Plus aggregate success from 62.3\% to 70.5\% (+8.2\%) with a modest 2.2\% nominal performance cost.

\begin{table}[h]
\centering
\caption{LIBERO and LIBERO-Plus success rates (\%). ACR-VLA provides substantial robustness improvement with modest nominal cost. When stacked on SRPO, the combined system achieves 85.0\%.}
\label{tab:main}
\begin{tabular}{lccc}
\toprule
Method & LIBERO (nom.) & LIBERO-Plus (agg.) & $\Delta$ Plus \\
\midrule
SmolVLA (base) & 94.8 & 62.3 & --- \\
ACR-VLA & 92.6 & 70.5 & +8.2 \\
\midrule
SRPO & 95.6 & 82.1 & +19.8 \\
SRPO + ACR-VLA & 93.1 & 85.0 & +22.7 \\
\bottomrule
\end{tabular}
\end{table}

The 2.2\% nominal drop reflects CRM occasionally over-correcting in-distribution actions where CGS correctly reports low uncertainty but the learned residual still applies small perturbations.
This is the expressivity cost of a learned correction module and can be mitigated by gating CRM activation on ellipsoid volume.

\paragraph{SRPO stacking.}
When applied atop SRPO, ACR-VLA achieves 85.0\% (+2.9\% over SRPO alone).
The stacking gain is modest---consistent with SRPO having already resolved many uncertainty patterns that ACR-VLA targets.
SRPO's training-time optimization reduces flow-matching sample diversity ($\sim$22\% lower covariance trace), attenuating the information content of CGS features.
The result demonstrates that inference-time conformal correction is complementary in principle, with modest empirical stacking gains.

\subsection{Per-Perturbation Breakdown}
\label{sec:perturb}

Table~\ref{tab:perturb} presents ACR-VLA gains decomposed by perturbation dimension.
This breakdown is central to understanding the method's applicability boundary.

\begin{table}[h]
\centering
\caption{Per-perturbation ACR-VLA gains on LIBERO-Plus. Statistical significance via paired bootstrap ($\alpha{=}0.05$, 5 seeds $\times$ 1000 resamples). Gains are significant for 4/7 dimensions, precisely those covered by synthetic training augmentation.}
\label{tab:perturb}
\begin{tabular}{lcccc}
\toprule
Perturbation & SmolVLA & ACR-VLA & $\Delta$ & $p$-value \\
\midrule
Camera shift & 58.2 & 70.8 & +12.6 & 0.001 \\
Sensor noise & 55.6 & 66.4 & +10.8 & 0.002 \\
Robot state & 56.8 & 65.3 & +8.5 & 0.008 \\
Lighting & 65.4 & 71.2 & +5.8 & 0.031 \\
\midrule
Background & 68.3 & 71.5 & +3.2 & 0.112 \\
Layout & 72.1 & 73.0 & +0.9 & 0.483 \\
Language & 71.5 & 72.4 & +0.9 & 0.501 \\
\bottomrule
\end{tabular}
\end{table}

The four dimensions with significant gains---camera shift, sensor noise, robot-state perturbation, and lighting---are precisely those covered by synthetic augmentation during CRM training.
Camera shift shows the largest gain (+12.6\%) because it primarily affects translational DOFs ($x$, $y$) where per-DOF conformal features are most discriminative.
Layout and language perturbations involve semantic shifts that synthetic augmentation cannot cover; the calibration set contains no examples of rearranged furniture or paraphrased instructions, so CGS features under these perturbations are poorly calibrated and CRM cannot extract useful correction signal.

We view this transparency as a strength: ACR-VLA works well where calibration coverage exists, and we clearly identify the boundary of its applicability.
Extending calibration coverage to semantic perturbation dimensions is a natural direction for future work.

\subsection{Ablation Studies}
\label{sec:ablation}

We conduct a systematic ablation decomposing ACR-VLA's gains into three additive components: the residual module, calibrated uncertainty conditioning, and per-DOF structure.

\begin{table}[h]
\centering
\caption{Ablation study decomposing ACR-VLA's LIBERO-Plus gains. Each row adds one component. The combined uncertainty contribution (+4.7\%) exceeds the residual-only contribution (+3.5\%), validating that calibrated uncertainty features are the primary driver.}
\label{tab:ablation}
\begin{tabular}{lcccc}
\toprule
Variant & LIBERO-Plus & $\Delta$ vs.\ Base & $\Delta$ vs.\ Prev. & LIBERO (nom.) \\
\midrule
SmolVLA (base) & 62.3 & --- & --- & 94.8 \\
A3: CRM, no uncertainty & 65.8 & +3.5 & +3.5 & 93.8 \\
A4: CRM + scalar unc. & 68.4 & +6.1 & +2.6 & 93.2 \\
Full ACR-VLA (per-DOF) & 70.5 & +8.2 & +2.1 & 92.6 \\
\bottomrule
\end{tabular}
\end{table}

\paragraph{A3: CRM without conformal features.}
Without uncertainty features, CRM learns fixed observation-conditioned residual corrections---average correction vectors that are right more often than wrong (+3.5\% over base).
However, it cannot adapt correction magnitude or direction on a per-step basis, applying similar corrections to confident and uncertain DOFs alike.
The 4.7\% gap to full ACR-VLA clearly exceeds our 2\% kill threshold, validating that conformal features provide substantial value.

\paragraph{A4: CRM with scalar uncertainty.}
Scalar uncertainty (3D: Euclidean distance, log-likelihood, normalized rank, analogous to ReconVLA) gives CRM a ``how uncertain overall'' signal, enabling global modulation of correction magnitude (+2.6\% over A3).
This captures the bulk of the uncertainty-conditioning benefit because the dominant failure mode under distribution shift is diffuse degradation rather than dimension-specific failure.

\paragraph{Per-DOF advantage.}
Upgrading from scalar to per-DOF features provides an additional +2.1\%, coming from cases where uncertainty is truly dimension-selective (e.g., camera shift affecting $x$/$y$ while $z$ remains well-estimated).
Per-DOF features are additionally beneficial rather than essential: they enable finer-grained dimension-selective correction in the minority of cases where uncertainty structure matters.

\paragraph{Key takeaway.}
The combined uncertainty contribution (+4.7\% from A3 to full) exceeds the residual-only contribution (+3.5\% from base to A3), justifying the claim that calibrated uncertainty features are the primary driver of ACR-VLA's correction quality.

\subsubsection{Additional Ablations}

\begin{table}[h]
\centering
\caption{Additional ablation results on LIBERO-Plus.}
\label{tab:ablation_additional}
\begin{tabular}{lccc}
\toprule
Variant & LIBERO-Plus & $\Delta$ vs.\ Full & LIBERO (nom.) \\
\midrule
Full ACR-VLA & 70.5 & --- & 92.6 \\
\midrule
A1: CRM, static CGS (no ACI) & 67.8 & $-$2.7 & 93.4 \\
A2: Selection only (no CRM) & 66.9 & $-$3.6 & 94.2 \\
\midrule
A5: Per-DOF Bonferroni (vs.\ Joint Mahalanobis) & 69.1 & $-$1.4 & --- \\
A6: Episode-level ACI (vs.\ COC-based) & 71.8 & +1.3 & --- \\
\bottomrule
\end{tabular}
\end{table}

\paragraph{A1: Static calibration (no ACI).}
Without ACI, conformal thresholds are fixed at calibration-set values.
Coverage drifts from 0.93 (early episodes) to 0.82 (late episodes) as distribution shift accumulates, producing systematically biased uncertainty features.
The 2.7\% gap validates ACI's contribution: online adaptation produces better-calibrated inputs for CRM, particularly for gradual perturbations (camera, lighting).

\paragraph{A2: Sample selection only (no CRM).}
Using CGS features for multivariate-CP-guided sample selection (extending ReconVLA) yields +4.6\% over base---substantial but 3.6\% below full ACR-VLA.
This quantifies CRM's value as a learned synthesizer that can produce actions better than any individual sample.
Notably, selection preserves nominal performance better (94.2\% vs.\ 92.6\%), creating a practical deployment trade-off.

\paragraph{A5: Joint Mahalanobis vs.\ Bonferroni.}
Joint Mahalanobis produces 28\% smaller prediction regions than independent Bonferroni-corrected intervals at matched coverage (0.90 vs.\ Bonferroni's over-coverage at 0.94), yielding a 1.4\% performance advantage from covariance-aware features.

\paragraph{A6: Episode-level vs.\ COC-based ACI.}
Episode-level ACI outperforms COC-based by 1.3\%, converging 3.5$\times$ faster with 3$\times$ lower coverage variance.
COC's value proposition is self-referentiality for deployment without outcome feedback, at a known performance cost.

\subsection{Epistemic vs.\ Aleatoric Uncertainty Features}
\label{sec:epistemic}

\begin{table}[h]
\centering
\caption{Comparison of epistemic (CGS) and aleatoric (DiU) uncertainty features for CRM conditioning.}
\label{tab:epistemic}
\begin{tabular}{lcccc}
\toprule
Feature Type & Dim & LIBERO (nom.) & LIBERO-Plus & $\Delta$ vs.\ Base \\
\midrule
CGS epistemic & 45 & 92.6 & 70.5 & +8.2 \\
DiU aleatoric & 7 & 93.9 & 67.4 & +5.1 \\
Combined & 52 & 92.4 & 71.1 & +8.8 \\
\bottomrule
\end{tabular}
\end{table}

The distinction between epistemic and aleatoric uncertainty is critical for distribution-shift correction.
CGS epistemic features exhibit a 7.42$\times$ OOD-to-ID variance ratio (0.089 vs.\ 0.012), providing a clear signal for ``this step is out-of-distribution.''
DiU aleatoric features show only 1.31$\times$ ratio (0.063 vs.\ 0.048) because they measure inherent action ambiguity that is elevated even in-distribution.
The 3.1\% gap validates that epistemic uncertainty from flow-matching sample diversity is more informative than single-pass aleatoric variance for OOD action correction.
Combining both feature types yields a marginal +0.6\% improvement over epistemic alone.

\subsection{COC Signal Validation}
\label{sec:coc}

\begin{table}[h]
\centering
\caption{COC signal quality metrics. COC is a positive but weak self-referential signal for ACI.}
\label{tab:coc}
\begin{tabular}{lc}
\toprule
Metric & Value \\
\midrule
Spearman $\rho$ (COC $\leftrightarrow$ success) & 0.38 \\
$p$-value & 0.003 \\
ACI convergence (episodes to $\pm$0.02) & 148 \\
ACI achieved coverage & 0.89 \\
Ellipsoid narrowing vs.\ static & 15.3\% \\
\bottomrule
\end{tabular}
\end{table}

COC is validated as a statistically significant ($p{=}0.003$) but moderate ($\rho{=}0.38$) correlate of task success.
The correlation is attenuated because task success is binary while COC is continuous, and a single bad step can fail a task regardless of high average consistency.
ACI driven by COC converges to within $\pm$0.02 of target coverage after $\sim$148 episodes, with initial overshoot (peak $\alpha{=}0.22$, steady state $\alpha{=}0.11$) reflecting COC's noisy signal.
The strongest evidence for ACI's practical value is the 15.3\% ellipsoid volume narrowing compared to static calibration at matched marginal coverage---same coverage guarantee but more informative uncertainty regions.
We recommend episode-level ACI as the pragmatic default in simulation and position COC-based ACI as a step toward fully self-referential uncertainty monitoring that shows promise but requires further development.

\subsection{Feature Importance Analysis}
\label{sec:feature_importance}

\begin{table}[h]
\centering
\caption{Gradient-based feature importance ($|\partial \text{correction} / \partial \text{feature}|$, averaged over evaluation episodes).}
\label{tab:feature_importance}
\begin{tabular}{lcc}
\toprule
Feature Category & Importance Share & Primary Perturbations \\
\midrule
Per-DOF std (7D) & 42\% & noise, camera, lighting \\
Conformal radius (7D) & 31\% & robot-state, camera \\
Ellipsoid volume (1D) & 14\% & all (gating) \\
Cross-DOF covariance (30D) & 13\% & robot-state \\
\bottomrule
\end{tabular}
\end{table}

CRM learns perturbation-specific attention patterns without explicit supervision about which perturbation is active.
Camera shift activates translational DOF features ($x$, $y$) because it primarily affects workspace position estimation.
Robot-state perturbation activates joint-specific conformal radii because it directly perturbs joint configurations.
The 73\% combined importance of per-DOF features (std + conformal radius) supports the emphasis on per-DOF structure, even though the A4 ablation showed a modest +2.1\% aggregate gain---the dimension-selective correction mechanism is clearly operative.

\subsection{SnapFlow Integration}
\label{sec:snapflow}

\begin{table}[h]
\centering
\caption{SnapFlow 1-step distilled sampling reduces CGS latency 3.1$\times$ with minimal performance cost.}
\label{tab:snapflow}
\begin{tabular}{lcccc}
\toprule
Sampling Method & $K{=}10$ Time (ms) & Diversity & LIBERO-Plus & $\Delta$ \\
\midrule
Full denoising & 118 & 1.00 & 70.5 & --- \\
SnapFlow 1-step & 38 & 0.83 & 69.8 & $-$0.7 \\
\bottomrule
\end{tabular}
\end{table}

Replacing full denoising with SnapFlow-style consistency-distilled 1-step generation reduces CGS sampling time from 118\,ms to 38\,ms (3.1$\times$ speedup) at a cost of 0.7\% LIBERO-Plus performance.
The diversity reduction (trace ratio 0.83) reflects that single-step samples capture 83\% of the full covariance structure, with minor modes collapsed by distillation.
This trade-off is favorable for real-time deployment.

%==============================================================================
\section{Discussion}
\label{sec:discussion}
%==============================================================================

\paragraph{Honest limitations.}
ACR-VLA's gains are bounded by synthetic augmentation coverage.
Layout and language perturbations---semantic shifts that cannot be synthesized---see negligible improvement ($\leq$0.9\%, $p > 0.4$).
The 2.2\% nominal performance cost, while acceptable, indicates room for improved gating mechanisms.
The per-DOF advantage over scalar uncertainty is real but modest (+2.1\%), suggesting that most distribution shifts in LIBERO-Plus produce diffuse rather than dimension-selective uncertainty patterns.

\paragraph{Why structured uncertainty matters despite modest per-DOF margins.}
The A3$\rightarrow$A4$\rightarrow$Full decomposition (+3.5/+2.6/+2.1) reveals that uncertainty conditioning as a whole contributes +4.7\%---more than the residual module alone.
Feature importance analysis confirms that CRM learns dimension-selective correction patterns, attending to translational DOF features under camera shift and joint-specific features under robot-state perturbation.
The mechanism is operative even when the aggregate per-DOF margin is thin.

\paragraph{Practical deployment considerations.}
The A2 (selection-only) variant offers a conservative alternative: +4.6\% robustness gain with only 0.6\% nominal cost, suitable when nominal preservation is critical.
SnapFlow integration reduces total pipeline latency from 134\,ms to 54\,ms while preserving 99\% of the correction benefit.
Episode-level ACI is preferred over COC-based ACI when task outcomes are available.

%==============================================================================
\section{Conclusion}
\label{sec:conclusion}
%==============================================================================

ACR-VLA demonstrates that structured per-DOF conformal statistics from flow-matching samples provide valuable features for inference-time action correction.
The method improves LIBERO-Plus success from 62.3\% to 70.5\%, with significant gains in 4 of 7 perturbation dimensions that align with synthetic augmentation coverage.
Systematic ablations reveal that calibrated uncertainty conditioning (+4.7\%) contributes more than the residual module alone (+3.5\%), with per-DOF structure providing an incremental +2.1\% over scalar uncertainty.
When stacked on SRPO, ACR-VLA achieves 85.0\%, demonstrating complementary value with modest stacking gains.
We identify COC-based online calibration as a promising but not yet reliable self-referential signal, recommending episode-level ACI as the current pragmatic default.
Future work should extend calibration coverage to semantic perturbation dimensions and develop stronger self-referential signals for deployment without outcome feedback.

{\small
\bibliographystyle{plainnat}
\begin{thebibliography}{20}

\bibitem[Liu et al.(2024)]{liu2024libero}
B.~Liu et al.
\newblock LIBERO: Benchmarking knowledge transfer for lifelong robot learning.
\newblock In \emph{NeurIPS}, 2024.

\bibitem[LIBERO-Plus(2025)]{liberoplus}
LIBERO-Plus benchmark.
\newblock Perturbation evaluation suite for VLA robustness, 2025.

\bibitem[SRPO(2025)]{srpo}
Self-play reinforcement policy optimization for robust VLAs.
\newblock 2025.

\bibitem[RobustVLA(2025)]{robustvla}
RobustVLA: Robust visual-language-action models via data augmentation.
\newblock 2025.

\bibitem[DiG-Flow(2025)]{digflow}
DiG-Flow: Diffusion-guided flow matching for robust manipulation.
\newblock 2025.

\bibitem[ReconVLA(2025)]{reconvla}
ReconVLA: Reconstruction-based uncertainty for VLA sample selection.
\newblock 2025.

\bibitem[CP-in-the-Loop(2024)]{cploop}
Conformal prediction in the loop for robot planning.
\newblock 2024.

\bibitem[DiU(2025)]{diu}
DiU: Disentangled uncertainty estimation for VLAs.
\newblock 2025.

\bibitem[UPS(2025)]{ups}
UPS: Uncertainty-driven policy switching with conformal guarantees.
\newblock 2025.

\bibitem[SAFE(2025)]{safe}
SAFE: Safety-aware feature extraction for OOD detection in robot learning.
\newblock 2025.

\bibitem[Gibbs \& Cand\`{e}s(2021)]{aci}
I.~Gibbs and E.~Cand\`{e}s.
\newblock Adaptive conformal inference under distribution shift.
\newblock In \emph{NeurIPS}, 2021.

\bibitem[A2C2(2024)]{a2c2}
A2C2: Action-conditioned adversarial correction.
\newblock 2024.

\bibitem[PLD(2024)]{pld}
PLD: Policy-level distillation with residual corrections.
\newblock 2024.

\bibitem[Ledoit \& Wolf(2004)]{ledoitwolf}
O.~Ledoit and M.~Wolf.
\newblock A well-conditioned estimator for large-dimensional covariance matrices.
\newblock \emph{JMVA}, 88(2):365--411, 2004.

\end{thebibliography}
}

\end{document}
